\documentclass[preview]{standalone}

\usepackage{amsmath}
\usepackage{amssymb}
\usepackage{stellar}
\usepackage{definitions}
\usepackage{bettelini}

\begin{document}

\id{linearalgebra-adjoints-spectral-theorem}
\genpage

\section{Riesz representation}

\begin{snippet}{riesz-vector-functional-map}
    Let \((V,\langle\cdot,\cdot\rangle)\) be an \innerproductspace over
    \(\mathbb{F}\in\{\realnumbers,\complexnumbers\}\). For every \(u\in V\),
    the map
    \[
        \rho_u\colon V\fromto\mathbb{F},
        \qquad
        \rho_u(v)=\langle v,u\rangle
    \]
    is a \linearfunctional because the \innerproduct is linear in the first
    variable.
\end{snippet}

\begin{snippettheorem}{riesz-representation-theorem}{Riesz representation theorem}
    Let \((V,\langle\cdot,\cdot\rangle)\) be a finite-dimensional
    \innerproductspace over
    \(\mathbb{F}\in\{\realnumbers,\complexnumbers\}\), and let
    \(\rho\in V^\ast\) be a \linearfunctional. There exists a unique
    \(u\in V\) such that
    \[
        \rho(v)=\langle v,u\rangle
    \]
    for every \(v\in V\).
\end{snippettheorem}

\begin{snippetproof}{riesz-representation-theorem-proof}{riesz-representation-theorem}{Riesz representation theorem}
    Let
    \[
        \mathcal{B}=\{e_1,\ldots,e_n\}
    \]
    be an \orthonormalbasis of \(V\). Define
    \[
        u=
        \overline{\rho(e_1)}e_1+\cdots+\overline{\rho(e_n)}e_n.
    \]
    If \(v\in V\), then by the \parsevalidentity,
    \[
        v=\langle v,e_1\rangle e_1+\cdots+\langle v,e_n\rangle e_n.
    \]
    Since \(\rho\) is linear,
    \[
        \rho(v)=
        \langle v,e_1\rangle\rho(e_1)
        +\cdots+
        \langle v,e_n\rangle\rho(e_n).
    \]
    On the other hand, using conjugate-linearity in the second variable,
    \[
        \langle v,u\rangle
        =
        \rho(e_1)\langle v,e_1\rangle
        +\cdots+
        \rho(e_n)\langle v,e_n\rangle.
    \]
    Hence \(\rho(v)=\langle v,u\rangle\) for every \(v\in V\).

    For uniqueness, suppose that \(u_1,u_2\in V\) both represent \(\rho\).
    Then
    \[
        \langle v,u_1-u_2\rangle=0
    \]
    for every \(v\in V\). Taking \(v=u_1-u_2\) gives
    \[
        \langle u_1-u_2,u_1-u_2\rangle=0.
    \]
    By definiteness of the \innerproduct, \(u_1-u_2=0\), so \(u_1=u_2\).
\end{snippetproof}

\section{Adjoint maps}

\begin{snippetdefinition}{adjoint-map-definition}{Adjoint map}
    Let \((V,\langle\cdot,\cdot\rangle_V)\) and
    \((W,\langle\cdot,\cdot\rangle_W)\) be finite-dimensional
    \innerproductspace[inner product spaces] over the same
    \field \(\mathbb{F}\in\{\realnumbers,\complexnumbers\}\), and let
    \(f\in\setoflineartransformations(V,W)\). The \emph{adjoint} of \(f\)
    is the unique map
    \[
        f^\ast\colon W\fromto V
    \]
    such that
    \[
        \langle v,f^\ast(w)\rangle_V
        =
        \langle f(v),w\rangle_W
    \]
    for every \(v\in V\) and every \(w\in W\).
\end{snippetdefinition}

\begin{snippetproof}{adjoint-map-definition-proof}{adjoint-map-definition}{Adjoint map}
    For each fixed \(w\in W\), the map
    \[
        V\fromto\mathbb{F},
        \qquad
        v\mapsto \langle f(v),w\rangle_W
    \]
    is a \linearfunctional on \(V\). By the \rieszrepresentation, there is
    a unique vector \(f^\ast(w)\in V\) such that
    \[
        \langle v,f^\ast(w)\rangle_V=\langle f(v),w\rangle_W
    \]
    for every \(v\in V\). This defines the adjoint map.
\end{snippetproof}

\begin{snippetproposition}{adjoint-map-linearity}{The adjoint is linear}
    Let \((V,\langle\cdot,\cdot\rangle_V)\) and
    \((W,\langle\cdot,\cdot\rangle_W)\) be finite-dimensional
    \innerproductspace[inner product spaces] over the same
    \field \(\mathbb{F}\in\{\realnumbers,\complexnumbers\}\), and let
    \(f\in\setoflineartransformations(V,W)\). Then
    \[
        f^\ast\in\setoflineartransformations(W,V).
    \]
\end{snippetproposition}

\begin{snippetproof}{adjoint-map-linearity-proof}{adjoint-map-linearity}{The adjoint is linear}
    Let \(w_1,w_2\in W\). For every \(v\in V\),
    \[
        \langle v,f^\ast(w_1+w_2)\rangle_V
        =
        \langle f(v),w_1+w_2\rangle_W
        =
        \langle v,f^\ast(w_1)+f^\ast(w_2)\rangle_V.
    \]
    By uniqueness in the \rieszrepresentation,
    \[
        f^\ast(w_1+w_2)=f^\ast(w_1)+f^\ast(w_2).
    \]
    If \(\lambda\in\mathbb{F}\) and \(w\in W\), then for every \(v\in V\),
    \[
        \langle v,f^\ast(\lambda w)\rangle_V
        =
        \langle f(v),\lambda w\rangle_W
        =
        \overline{\lambda}\langle f(v),w\rangle_W
        =
        \langle v,\lambda f^\ast(w)\rangle_V.
    \]
    Again by uniqueness,
    \[
        f^\ast(\lambda w)=\lambda f^\ast(w).
    \]
    Thus \(f^\ast\) is linear.
\end{snippetproof}

\begin{snippetexample}{adjoint-euclidean-map-example}{Adjoint of a map between Euclidean spaces}
    Let
    \[
        f\colon\realnumbers^3\fromto\realnumbers^2,
        \qquad
        f(x,y,z)=(y+3z,2x),
    \]
    and use the canonical \innerproduct[inner products]. The \adjointmap of
    \(f\) is
    \[
        f^\ast\colon\realnumbers^2\fromto\realnumbers^3,
        \qquad
        f^\ast(s,t)=(2t,s,3s).
    \]
    Indeed, for every \((x,y,z)\in\realnumbers^3\) and
    \((s,t)\in\realnumbers^2\),
    \[
        \langle f(x,y,z),(s,t)\rangle
        =
        (y+3z)s+2xt
        =
        \langle (x,y,z),(2t,s,3s)\rangle.
    \]
\end{snippetexample}

\begin{snippetexample}{rank-one-map-adjoint-example}{Adjoint of a rank-one map}
    Let \((V,\langle\cdot,\cdot\rangle_V)\) and
    \((W,\langle\cdot,\cdot\rangle_W)\) be finite-dimensional
    \innerproductspace[inner product spaces] over
    \(\mathbb{F}\in\{\realnumbers,\complexnumbers\}\). Fix \(x\in V\) and
    \(w\in W\), and define
    \[
        f\colon V\fromto W,
        \qquad
        f(v)=\langle v,x\rangle_V\,w.
    \]
    Then
    \[
        f^\ast(y)=\langle y,w\rangle_W\,x.
    \]
    Indeed, for every \(v\in V\) and \(y\in W\),
    \[
        \langle f(v),y\rangle_W
        =
        \langle v,x\rangle_V\langle w,y\rangle_W
        =
        \langle v,\langle y,w\rangle_W x\rangle_V.
    \]
\end{snippetexample}

\begin{snippetproposition}{adjoint-formal-properties}{Formal properties of the adjoint}
    Let \(V,W,U\) be finite-dimensional \innerproductspace[inner product
    spaces] over \(\mathbb{F}\in\{\realnumbers,\complexnumbers\}\). If
    \(f,g\in\setoflineartransformations(V,W)\),
    \(h\in\setoflineartransformations(W,U)\), and
    \(\lambda\in\mathbb{F}\), then:
    \begin{enumerate}
        \item \((f+g)^\ast=f^\ast+g^\ast\);
        \item \((\lambda f)^\ast=\overline{\lambda}f^\ast\);
        \item \((f^\ast)^\ast=f\);
        \item \((\identityfunc_V)^\ast=\identityfunc_V\);
        \item \((h\circ f)^\ast=f^\ast\circ h^\ast\).
    \end{enumerate}
\end{snippetproposition}

\begin{snippetproof}{adjoint-formal-properties-proof}{adjoint-formal-properties}{Formal properties of the adjoint}
    The first two identities follow from the defining property and uniqueness
    of the \adjointmap:
    \[
        \langle v,(f+g)^\ast(w)\rangle
        =
        \langle (f+g)(v),w\rangle
        =
        \langle v,f^\ast(w)+g^\ast(w)\rangle
    \]
    and
    \[
        \langle v,(\lambda f)^\ast(w)\rangle
        =
        \lambda\langle f(v),w\rangle
        =
        \langle v,\overline{\lambda}f^\ast(w)\rangle.
    \]
    For the double adjoint, for every \(v\in V\) and \(w\in W\),
    \[
        \langle w,f(v)\rangle_W
        =
        \overline{\langle f(v),w\rangle_W}
        =
        \overline{\langle v,f^\ast(w)\rangle_V}
        =
        \langle f^\ast(w),v\rangle_V.
    \]
    Thus \((f^\ast)^\ast=f\). The identity map satisfies
    \[
        \langle v,\identityfunc_V(w)\rangle
        =
        \langle \identityfunc_V(v),w\rangle,
    \]
    so \((\identityfunc_V)^\ast=\identityfunc_V\).

    Finally, for \(v\in V\) and \(u\in U\),
    \[
        \langle v,(h\circ f)^\ast(u)\rangle_V
        =
        \langle h(f(v)),u\rangle_U
        =
        \langle v,f^\ast(h^\ast(u))\rangle_V.
    \]
    Hence \((h\circ f)^\ast=f^\ast\circ h^\ast\).
\end{snippetproof}

\section{Matrix of the adjoint}

\includesnpt{matrix-conjugate-transpose-definition}

\begin{snippetproposition}{adjoint-matrix-orthonormal-bases}{Matrix of the adjoint in orthonormal bases}
    Let \((V,\langle\cdot,\cdot\rangle_V)\) and
    \((W,\langle\cdot,\cdot\rangle_W)\) be finite-dimensional
    \innerproductspace[inner product spaces] over
    \(\mathbb{F}\in\{\realnumbers,\complexnumbers\}\). Let
    \(\mathcal{B}=\{e_1,\ldots,e_n\}\) be an \orthonormalbasis of \(V\),
    and let \(\mathcal{C}=\{w_1,\ldots,w_m\}\) be an
    \orthonormalbasis of \(W\). If
    \(f\in\setoflineartransformations(V,W)\), then
    \[
        M(f^\ast,\mathcal{C},\mathcal{B})
        =
        M(f,\mathcal{B},\mathcal{C})^\conjtrans.
    \]
\end{snippetproposition}

\begin{snippetproof}{adjoint-matrix-orthonormal-bases-proof}{adjoint-matrix-orthonormal-bases}{Matrix of the adjoint in orthonormal bases}
    Put
    \[
        A=M(f,\mathcal{B},\mathcal{C})=(a_{jk}),
        \qquad
        B=M(f^\ast,\mathcal{C},\mathcal{B})=(b_{kj}).
    \]
    Since \(\mathcal{C}\) is an \orthonormalbasis, the \(j\)-th coordinate
    of \(f(e_k)\) is
    \[
        a_{jk}=\langle f(e_k),w_j\rangle_W.
    \]
    Since \(\mathcal{B}\) is an \orthonormalbasis, the \(k\)-th coordinate
    of \(f^\ast(w_j)\) is
    \[
        b_{kj}=\langle f^\ast(w_j),e_k\rangle_V.
    \]
    By the definition of \adjointmap,
    \[
        \langle e_k,f^\ast(w_j)\rangle_V
        =
        \langle f(e_k),w_j\rangle_W.
    \]
    Therefore
    \[
        b_{kj}
        =
        \overline{\langle e_k,f^\ast(w_j)\rangle_V}
        =
        \overline{\langle f(e_k),w_j\rangle_W}
        =
        \overline{a_{jk}}.
    \]
    Hence \(B=A^\conjtrans\).
\end{snippetproof}

\section{Self-adjoint endomorphisms}

\begin{snippetdefinition}{self-adjoint-endomorphism-definition}{Self-adjoint endomorphism}
    Let \((V,\langle\cdot,\cdot\rangle)\) be a finite-dimensional
    \innerproductspace, and let
    \(f\in\setoflineartransformations(V,V)\). The endomorphism \(f\) is
    \emph{self-adjoint} if
    \[
        f^\ast=f.
    \]
    Equivalently,
    \[
        \langle f(u),v\rangle=\langle u,f(v)\rangle
    \]
    for every \(u,v\in V\).
\end{snippetdefinition}

\begin{snippetproposition}{self-adjoint-matrix-characterization}{Matrices of self-adjoint endomorphisms}
    Let \((V,\langle\cdot,\cdot\rangle)\) be a finite-dimensional
    \innerproductspace over
    \(\mathbb{F}\in\{\realnumbers,\complexnumbers\}\), let
    \(f\in\setoflineartransformations(V,V)\), and let \(\mathcal{B}\) be an
    \orthonormalbasis of \(V\). Then \(f\) is \selfadjoint if and only if
    \[
        M(f,\mathcal{B},\mathcal{B})^\conjtrans
        =
        M(f,\mathcal{B},\mathcal{B}).
    \]
    In the real case, this means that the matrix is a \symmetricmatrix. In
    the complex case, this means that the matrix is \hermitianmatrix[Hermitian].
\end{snippetproposition}

\begin{snippetproof}{self-adjoint-matrix-characterization-proof}{self-adjoint-matrix-characterization}{Matrices of self-adjoint endomorphisms}
    Since \(\mathcal{B}\) is an \orthonormalbasis,
    \[
        M(f^\ast,\mathcal{B},\mathcal{B})
        =
        M(f,\mathcal{B},\mathcal{B})^\conjtrans.
    \]
    Therefore
    \[
        f=f^\ast
        \Longleftrightarrow
        M(f,\mathcal{B},\mathcal{B})
        =
        M(f,\mathcal{B},\mathcal{B})^\conjtrans.
    \]
    If \(\mathbb{F}=\realnumbers\), the conjugate transpose is the
    transpose. If \(\mathbb{F}=\complexnumbers\), the condition is exactly
    the \hermitianmatrix[Hermitian] condition.
\end{snippetproof}

\begin{snippetnote}{self-adjoint-matrix-needs-orthonormal-basis}{Self-adjoint matrices require orthonormal bases}
    Let \(\realnumbers^2\) have the canonical \innerproduct, and let
    \[
        f(x,y)=(x,2y).
    \]
    Then \(f\) is \selfadjoint, because its matrix in the canonical
    \orthonormalbasis is
    \[
        \begin{pmatrix}
            1 & 0\\
            0 & 2
        \end{pmatrix}.
    \]
    However, with respect to the non-orthonormal \basis
    \[
        \mathcal{B}=\{(1,1),(1,2)\},
    \]
    the associated matrix is
    \[
        M(f,\mathcal{B},\mathcal{B})
        =
        \begin{pmatrix}
            0 & -2\\
            1 & 3
        \end{pmatrix},
    \]
    which is not \symmetricmatrix[symmetric].
\end{snippetnote}

\section{Normal endomorphisms}

\begin{snippetdefinition}{normal-endomorphism-definition}{Normal endomorphism}
    Let \((V,\langle\cdot,\cdot\rangle)\) be a finite-dimensional
    \innerproductspace, and let
    \(f\in\setoflineartransformations(V,V)\). The endomorphism \(f\) is
    \emph{normal} if
    \[
        \|f(v)\|=\|f^\ast(v)\|
    \]
    for every \(v\in V\).
\end{snippetdefinition}

\begin{snippetproposition}{self-adjoint-endomorphisms-are-normal}{Self-adjoint endomorphisms are normal}
    Let \((V,\langle\cdot,\cdot\rangle)\) be a finite-dimensional
    \innerproductspace. Every \selfadjoint endomorphism of \(V\) is a
    \normalendomorphism.
\end{snippetproposition}

\begin{snippetproof}{self-adjoint-endomorphisms-are-normal-proof}{self-adjoint-endomorphisms-are-normal}{Self-adjoint endomorphisms are normal}
    If \(f=f^\ast\), then
    \[
        \|f(v)\|=\|f^\ast(v)\|
    \]
    for every \(v\in V\). Hence \(f\) is a \normalendomorphism.
\end{snippetproof}

\begin{snippetproposition}{normal-endomorphism-characterizations}{Characterizations of normal endomorphisms}
    Let \((V,\langle\cdot,\cdot\rangle)\) be a finite-dimensional
    \innerproductspace over
    \(\mathbb{F}\in\{\realnumbers,\complexnumbers\}\), and let
    \(f\in\setoflineartransformations(V,V)\). The following are equivalent:
    \begin{enumerate}
        \item \(f\) is a \normalendomorphism;
        \item \(\langle f(v),f(w)\rangle=\langle f^\ast(v),f^\ast(w)\rangle\) for every \(v,w\in V\);
        \item \(f\circ f^\ast=f^\ast\circ f\).
    \end{enumerate}
\end{snippetproposition}

\begin{snippetproof}{normal-endomorphism-characterizations-proof}{normal-endomorphism-characterizations}{Characterizations of normal endomorphisms}
    The implication \(2\Rightarrow1\) follows by taking \(w=v\).

    Suppose \(1\) holds. If \(\mathbb{F}=\realnumbers\), use the real
    polarization identity
    \[
        \langle x,y\rangle
        =
        \frac{\|x+y\|^2-\|x-y\|^2}{4}.
    \]
    If \(\mathbb{F}=\complexnumbers\), use the complex polarization identity
    \[
        \langle x,y\rangle
        =
        \frac{\|x+y\|^2-\|x-y\|^2}{4}
        +
        \frac{\mathrm{i}\|x+\mathrm{i}y\|^2-\mathrm{i}\|x-\mathrm{i}y\|^2}{4}.
    \]
    Since normality gives \(\|f(z)\|=\|f^\ast(z)\|\) for every \(z\in V\),
    applying the appropriate polarization identity gives
    \[
        \langle f(v),f(w)\rangle
        =
        \langle f^\ast(v),f^\ast(w)\rangle.
    \]
    Thus \(1\Rightarrow2\).

    For every \(v,w\in V\),
    \[
        \langle f(v),f(w)\rangle
        =
        \langle v,(f^\ast\circ f)(w)\rangle
    \]
    and
    \[
        \langle f^\ast(v),f^\ast(w)\rangle
        =
        \langle v,(f\circ f^\ast)(w)\rangle.
    \]
    Hence condition \(2\) holds if and only if
    \[
        \langle v,(f^\ast\circ f)(w)-(f\circ f^\ast)(w)\rangle=0
    \]
    for every \(v,w\in V\), which by definiteness is equivalent to
    \[
        f^\ast\circ f=f\circ f^\ast.
    \]
\end{snippetproof}

\begin{snippetexample}{normal-not-self-adjoint-example}{A normal endomorphism that is not self-adjoint}
    Let \(A\in\matrices_{2\times2}(\complexnumbers)\) be
    \[
        A=
        \begin{pmatrix}
            1 & \mathrm{i}\\
            1 & 1
        \end{pmatrix}.
    \]
    Then
    \[
        A^\conjtrans=
        \begin{pmatrix}
            1 & 1\\
            -\mathrm{i} & 1
        \end{pmatrix}.
    \]
    Since \(A\neq A^\conjtrans\), the endomorphism \(L_A\) of
    \(\complexnumbers^2\) is not \selfadjoint. However,
    \[
        AA^\conjtrans
        =
        \begin{pmatrix}
            2 & 1+\mathrm{i}\\
            1-\mathrm{i} & 2
        \end{pmatrix}
        =
        A^\conjtrans A.
    \]
    Hence \(L_A\) is a \normalendomorphism.
\end{snippetexample}

\begin{snippetlemma}{normal-eigenvectors-of-adjoint}{Eigenvectors of a normal endomorphism and its adjoint}
    Let \((V,\langle\cdot,\cdot\rangle)\) be a finite-dimensional
    \innerproductspace over
    \(\mathbb{F}\in\{\realnumbers,\complexnumbers\}\), and let
    \(f\in\setoflineartransformations(V,V)\) be a \normalendomorphism. Then, for every
    \(\lambda\in\mathbb{F}\) and every \(v\in V\),
    \[
        f(v)=\lambda v
        \qquad
        \Longleftrightarrow
        \qquad
        f^\ast(v)=\overline{\lambda}v.
    \]
\end{snippetlemma}

\begin{snippetproof}{normal-eigenvectors-of-adjoint-proof}{normal-eigenvectors-of-adjoint}{Eigenvectors of a normal endomorphism and its adjoint}
    By the commutation characterization of normality, the endomorphism
    \[
        f-\lambda\identityfunc_V
    \]
    is a \normalendomorphism. Its adjoint is
    \[
        f^\ast-\overline{\lambda}\identityfunc_V.
    \]
    Hence, for every \(v\in V\),
    \[
        \|(f-\lambda\identityfunc_V)(v)\|
        =
        \|(f^\ast-\overline{\lambda}\identityfunc_V)(v)\|.
    \]
    The first norm is zero if and only if the second norm is zero. This is
    exactly the claimed equivalence.
\end{snippetproof}

\begin{snippetlemma}{normal-eigenvectors-distinct-eigenvalues-orthogonal}{Eigenvectors for distinct eigenvalues of a normal endomorphism are orthogonal}
    Let \((V,\langle\cdot,\cdot\rangle)\) be a finite-dimensional
    \innerproductspace over
    \(\mathbb{F}\in\{\realnumbers,\complexnumbers\}\), and let
    \(f\in\setoflineartransformations(V,V)\) be a \normalendomorphism. If
    \(v\in V_\lambda\), \(w\in V_\mu\), and \(\lambda\neq\mu\), then
    \[
        v\perp w.
    \]
\end{snippetlemma}

\begin{snippetproof}{normal-eigenvectors-distinct-eigenvalues-orthogonal-proof}{normal-eigenvectors-distinct-eigenvalues-orthogonal}{Eigenvectors for distinct eigenvalues of a normal endomorphism are orthogonal}
    Since \(w\in V_\mu\), the previous lemma gives
    \[
        f^\ast(w)=\overline{\mu}w.
    \]
    Therefore
    \[
        \lambda\langle v,w\rangle
        =
        \langle f(v),w\rangle
        =
        \langle v,f^\ast(w)\rangle
        =
        \langle v,\overline{\mu}w\rangle
        =
        \mu\langle v,w\rangle.
    \]
    Since \(\lambda\neq\mu\), it follows that \(\langle v,w\rangle=0\).
\end{snippetproof}

\section{Self-adjoint eigenvalues}

\begin{snippetlemma}{complex-self-adjoint-eigenvalues-real}{Eigenvalues of complex self-adjoint endomorphisms}
    Let \((V,\langle\cdot,\cdot\rangle)\) be a finite-dimensional complex
    \innerproductspace, and let \(f\in\setoflineartransformations(V,V)\) be
    \selfadjoint. Every \eigenvalue of \(f\) is real.
\end{snippetlemma}

\begin{snippetproof}{complex-self-adjoint-eigenvalues-real-proof}{complex-self-adjoint-eigenvalues-real}{Eigenvalues of complex self-adjoint endomorphisms}
    Let \(\lambda\in\complexnumbers\) be an \eigenvalue of \(f\), and let
    \(v\neq0\) satisfy \(f(v)=\lambda v\). Since \(f=f^\ast\),
    \[
        \lambda\|v\|^2
        =
        \langle f(v),v\rangle
        =
        \langle v,f(v)\rangle
        =
        \overline{\lambda}\|v\|^2.
    \]
    Since \(v\neq0\), \(\|v\|^2\neq0\), and therefore
    \(\lambda=\overline{\lambda}\). Thus \(\lambda\in\realnumbers\).
\end{snippetproof}

\begin{snippetlemma}{real-self-adjoint-characteristic-roots-real}{Complex roots of a real self-adjoint characteristic polynomial}
    Let \((V,\langle\cdot,\cdot\rangle)\) be a finite-dimensional real
    \innerproductspace, and let \(f\in\setoflineartransformations(V,V)\) be
    \selfadjoint. Every complex root of the \characteristicpolynomial of
    \(f\) is real.
\end{snippetlemma}

\begin{snippetproof}{real-self-adjoint-characteristic-roots-real-proof}{real-self-adjoint-characteristic-roots-real}{Complex roots of a real self-adjoint characteristic polynomial}
    Let \(\mathcal{B}\) be an \orthonormalbasis of \(V\), and put
    \[
        A=M(f,\mathcal{B},\mathcal{B})\in\matrices_{n\times n}(\realnumbers).
    \]
    Since \(f\) is \selfadjoint, \(A^\transpose=A\). Regard \(A\) as a
    complex matrix. Then \(A^\conjtrans=A\), so the induced endomorphism
    \(L_A\) of \(\complexnumbers^n\) is complex self-adjoint. By the complex
    case, all complex roots of its characteristic polynomial are real. This
    is the same polynomial as the characteristic polynomial of \(f\).
\end{snippetproof}

\section{Spectral theorem}

\begin{snippettheorem}{abstract-spectral-theorem}{Abstract spectral theorem}
    Let \((V,\langle\cdot,\cdot\rangle)\) be a finite-dimensional
    \innerproductspace over
    \(\mathbb{F}\in\{\realnumbers,\complexnumbers\}\), and let
    \(f\in\setoflineartransformations(V,V)\). Then \(V\) has an
    \orthonormalbasis of \eigenvector[eigenvectors] of \(f\) if and only if
    \begin{enumerate}
        \item \(f\) is a \normalendomorphism;
        \item the \characteristicpolynomial of \(f\) splits completely in \(\mathbb{F}[t]\).
    \end{enumerate}
\end{snippettheorem}

\begin{snippetproof}{abstract-spectral-theorem-proof}{abstract-spectral-theorem}{Abstract spectral theorem}
    Suppose that
    \[
        \mathcal{B}=\{e_1,\ldots,e_n\}
    \]
    is an \orthonormalbasis of \eigenvector[eigenvectors] of \(f\), with
    \(f(e_i)=\lambda_i e_i\). Then
    \[
        p_f(t)=(t-\lambda_1)\cdots(t-\lambda_n)
    \]
    splits in \(\mathbb{F}[t]\). Also, by the \parsevalidentity, every
    \(v\in V\) satisfies
    \[
        v=\langle v,e_1\rangle e_1+\cdots+\langle v,e_n\rangle e_n.
    \]
    Since \(f^\ast(e_i)=\overline{\lambda_i}e_i\),
    \[
        \|f(v)\|^2
        =
        \sum_{i=1}^{n}|\lambda_i|^2|\langle v,e_i\rangle|^2
        =
        \|f^\ast(v)\|^2.
    \]
    Hence \(f\) is a \normalendomorphism.

    Conversely, suppose that \(f\) is a \normalendomorphism and that \(p_f(t)\) splits in
    \(\mathbb{F}[t]\). We prove the statement by induction on
    \(n=\lineardim V\). If \(n=1\), any normalized nonzero vector is an
    \orthonormalbasis of eigenvectors.

    Let \(n>1\). Since \(p_f(t)\) splits, \(f\) has an eigenvalue
    \(\lambda\in\mathbb{F}\). Choose a normalized eigenvector \(v\) with
    \(f(v)=\lambda v\), and put
    \[
        W=\linearspan(v)^\perp
        =
        \{w\in V\suchthat \langle w,v\rangle=0\}.
    \]
    Then
    \[
        V=\linearspan(v)\oplus W
    \]
    and \(\lineardim W=n-1\). If \(w\in W\), then by normality
    \(f^\ast(v)=\overline{\lambda}v\), so
    \[
        \langle f(w),v\rangle
        =
        \langle w,f^\ast(v)\rangle
        =
        \lambda\langle w,v\rangle
        =
        0.
    \]
    Thus \(W\) is stable under \(f\). A similar computation using
    \(f(v)=\lambda v\) shows that \(W\) is stable under \(f^\ast\). Hence
    the adjoint of \(f_{\mid W}\) is \(f^\ast_{\mid W}\), and
    \(f_{\mid W}\) is a \normalendomorphism.

    With respect to the decomposition
    \(V=\linearspan(v)\oplus W\), the characteristic polynomial factors as
    \[
        p_f(t)=(t-\lambda)p_{f_{\mid W}}(t).
    \]
    Therefore \(p_{f_{\mid W}}(t)\) also splits in \(\mathbb{F}[t]\). By the
    induction hypothesis, \(W\) has an \orthonormalbasis
    \(\{v_2,\ldots,v_n\}\) of eigenvectors of \(f_{\mid W}\). Then
    \[
        \{v,v_2,\ldots,v_n\}
    \]
    is an \orthonormalbasis of \(V\) consisting of eigenvectors of \(f\).
\end{snippetproof}

\begin{snippetcorollary}{complex-spectral-theorem}{Complex spectral theorem}
    Let \((V,\langle\cdot,\cdot\rangle)\) be a finite-dimensional complex
    \innerproductspace, and let \(f\in\setoflineartransformations(V,V)\).
    Then \(V\) has an \orthonormalbasis of \eigenvector[eigenvectors] of
    \(f\) if and only if \(f\) is a \normalendomorphism.
\end{snippetcorollary}

\begin{snippetproof}{complex-spectral-theorem-proof}{complex-spectral-theorem}{Complex spectral theorem}
    By the fundamental theorem of algebra, every complex characteristic
    polynomial splits in \(\complexnumbers[t]\). The result follows from
    the abstract \spectraltheorem.
\end{snippetproof}

\begin{snippetcorollary}{real-self-adjoint-spectral-theorem}{Real spectral theorem for self-adjoint endomorphisms}
    Let \((V,\langle\cdot,\cdot\rangle)\) be a finite-dimensional real
    \innerproductspace, and let \(f\in\setoflineartransformations(V,V)\).
    Then \(V\) has an \orthonormalbasis of \eigenvector[eigenvectors] of
    \(f\) if and only if \(f\) is \selfadjoint.
\end{snippetcorollary}

\begin{snippetproof}{real-self-adjoint-spectral-theorem-proof}{real-self-adjoint-spectral-theorem}{Real spectral theorem for self-adjoint endomorphisms}
    If \(V\) has an \orthonormalbasis of eigenvectors of \(f\), then the
    associated matrix of \(f\) in that basis is a real diagonal matrix, hence
    a \symmetricmatrix. Since the basis is \orthonormalbasis[orthonormal], \(f\) is \selfadjoint.

    Conversely, if \(f\) is \selfadjoint, then \(f\) is a \normalendomorphism. Moreover,
    every complex root of \(p_f(t)\) is real, so \(p_f(t)\) splits in
    \(\realnumbers[t]\). The abstract \spectraltheorem gives an
    \orthonormalbasis of eigenvectors.
\end{snippetproof}

\section{Matrix spectral theorems}

\begin{snippetdefinition}{orthogonal-matrix-definition}{Orthogonal matrix}
    A real square \matrix \(P\in\matrices_{n\times n}(\realnumbers)\) is
    \emph{orthogonal} if
    \[
        P^\transpose P=\identmatrix{n}.
    \]
    Equivalently, \(P^{-1}=P^\transpose\).
\end{snippetdefinition}

\begin{snippetdefinition}{unitary-matrix-definition}{Unitary matrix}
    A complex square \matrix \(U\in\matrices_{n\times n}(\complexnumbers)\)
    is \emph{unitary} if
    \[
        U^\conjtrans U=\identmatrix{n}.
    \]
    Equivalently, \(U^{-1}=U^\conjtrans\).
\end{snippetdefinition}

\begin{snippetdefinition}{normal-matrix-definition}{Normal matrix}
    A complex square \matrix \(A\in\matrices_{n\times n}(\complexnumbers)\)
    is \emph{normal} if
    \[
        AA^\conjtrans=A^\conjtrans A.
    \]
\end{snippetdefinition}

\begin{snippettheorem}{real-symmetric-matrix-spectral-theorem}{Real spectral theorem for symmetric matrices}
    Let \(A\in\matrices_{n\times n}(\realnumbers)\). Then \(A\) is a
    \symmetricmatrix if and only if there exist an \orthogonalmatrix
    \(P\in\matrices_{n\times n}(\realnumbers)\) and a real diagonal matrix
    \(D\in\matrices_{n\times n}(\realnumbers)\) such that
    \[
        P^\transpose AP=D.
    \]
\end{snippettheorem}

\begin{snippetproof}{real-symmetric-matrix-spectral-theorem-proof}{real-symmetric-matrix-spectral-theorem}{Real spectral theorem for symmetric matrices}
    Let \(L_A\colon\realnumbers^n\fromto\realnumbers^n\) be the
    endomorphism \(L_A(x)=Ax\), with the canonical \innerproduct. The
    canonical basis is an \orthonormalbasis, and \(L_A\) is \selfadjoint if
    and only if \(A^\transpose=A\). By the real spectral theorem for
    self-adjoint endomorphisms, this is equivalent to the existence of an
    \orthonormalbasis of eigenvectors. Placing those eigenvectors as the
    columns of \(P\) gives an \orthogonalmatrix and a diagonal matrix \(D\)
    satisfying
    \[
        P^\transpose AP=D.
    \]
    Conversely, if such \(P\) and \(D\) exist, then
    \[
        A=PDP^\transpose,
    \]
    so \(A^\transpose=PD^\transpose P^\transpose=PDP^\transpose=A\).
\end{snippetproof}

\begin{snippettheorem}{complex-hermitian-matrix-spectral-theorem}{Complex spectral theorem for Hermitian matrices}
    Let \(A\in\matrices_{n\times n}(\complexnumbers)\). Then \(A\) is
    \hermitianmatrix[Hermitian] if and only if there exist a \unitarymatrix
    \(U\in\matrices_{n\times n}(\complexnumbers)\) and a real diagonal
    matrix \(D\in\matrices_{n\times n}(\realnumbers)\) such that
    \[
        U^\conjtrans AU=D.
    \]
\end{snippettheorem}

\begin{snippetproof}{complex-hermitian-matrix-spectral-theorem-proof}{complex-hermitian-matrix-spectral-theorem}{Complex spectral theorem for Hermitian matrices}
    Let \(L_A\colon\complexnumbers^n\fromto\complexnumbers^n\) be
    \(L_A(x)=Ax\), with the canonical \innerproduct. Then \(L_A\) is
    \selfadjoint if and only if \(A^\conjtrans=A\), that is, if and only if
    \(A\) is \hermitianmatrix[Hermitian]. A complex \selfadjoint endomorphism is a \normalendomorphism, so by
    the complex \spectraltheorem it has an \orthonormalbasis of
    eigenvectors. Its eigenvalues are real. Placing the eigenvectors as the
    columns of \(U\) gives a \unitarymatrix and a real diagonal matrix \(D\)
    such that
    \[
        U^\conjtrans AU=D.
    \]
    Conversely, if such \(U\) and real diagonal \(D\) exist, then
    \[
        A=UDU^\conjtrans,
    \]
    so
    \[
        A^\conjtrans=UD^\conjtrans U^\conjtrans=UDU^\conjtrans=A.
    \]
\end{snippetproof}

\begin{snippetcorollary}{complex-normal-matrix-spectral-theorem}{Complex spectral theorem for normal matrices}
    Let \(A\in\matrices_{n\times n}(\complexnumbers)\). Then \(A\) is a
    \normalmatrix if and only if there exist a \unitarymatrix
    \(U\in\matrices_{n\times n}(\complexnumbers)\) and a complex diagonal
    matrix \(D\in\matrices_{n\times n}(\complexnumbers)\) such that
    \[
        U^\conjtrans AU=D.
    \]
\end{snippetcorollary}

\begin{snippetproof}{complex-normal-matrix-spectral-theorem-proof}{complex-normal-matrix-spectral-theorem}{Complex spectral theorem for normal matrices}
    The matrix \(A\) is a \normalmatrix exactly when the endomorphism
    \(L_A\colon\complexnumbers^n\fromto\complexnumbers^n\) is a \normalendomorphism. By
    the complex \spectraltheorem, this is equivalent to the existence of an
    \orthonormalbasis of eigenvectors of \(L_A\). Placing those eigenvectors
    as the columns of \(U\) gives a \unitarymatrix and a diagonal matrix
    \(D\) satisfying
    \[
        U^\conjtrans AU=D.
    \]
\end{snippetproof}

\end{document}
